\begin{textsample}{POS Dim 1 – vem\_ed\_19 – Score 10.00 – vem\_ed\_19/t090.txt}  \label{ex:f1_pos_052}
\textbf{continuação} Reflexões: gestão dos projetos Em 13 de junho, Ana Carolina Freschi, do \textbf{apoio} aos PCIs em inglês na equipe PCI / Cesu, apresentou o trabalho com seus colegas: Divinia Jithoo, Lindelwa Mkhize (Durban University of Technology, África do Sul), Daniel Otieno Okech (Kenyatta University, Quênia), Eva Haug (AUAS, Holanda) e Harshita Tripati (Shiv Nadar University, Índia).
Em pauta, a dinâmica de facilitação de Intercâmbios Virtuais nos diversos \textbf{contextos} educacionais desses países — sob os \textbf{pontos} de vista tecnológico, intercultural, acadêmico e \textbf{linguístico}.
Nesse mesmo dia, o painel “¡ SOS!
Intervención del Coordinador COIL” \textbf{foi} conduzido por Succi Junior, Gisselle Morales Veloquio (diretora de aprendizagem global do Tecnológico de Monterrey, México) e Regiane Moreira (professora responsável pelo \textbf{apoio} aos PCIs em espanhol nas Fatecs).
Os autores sugerem 10 passos para a intervenção do coordenador de Intercâmbios Virtuais, organizados nas seguintes etapas: institucionalizar; definir claramente o \textbf{papel} do coordenador COIL; \textbf{alinhar} \textbf{expectativas} da instituição, seus coordenadores, professores e parceiros; determinar critérios para busca de parceiros; oferecer formação continuada (para professores novos e experientes em Intercâmbios Virtuais); gerenciar o design dos projetos; realizar o acompanhamento dos projetos; \textbf{criar} índices de comparação e de informação; explicar e difundir constantemente os projetos COIL; decidir quando intervir (intervenção de emergência, apenas quando necessário).

% matched lemmas: alinhar, apoio, contexto, continuação, criar, expectativa, linguístico, papel, ponto, ser
\end{textsample}
